\subsection*{21. Homogeneity Is Valuable, but Too Much Specificity Produces a Mirage}

Define a segment so narrowly that every buyer has the same job, organization size, technology stack, city, seniority, purchasing date, budget, and personality, and you may find 3 perfect customers. The proposition will seem precise precisely because there is no meaningful market left.

A good segment balances \textbf{within-group similarity} against \textbf{economic density.} Its members should resemble one another on the handful of variables that cause the relevant purchase and make delivery repeatable. The remaining differences can be tolerated or handled by ordinary product variation.

A quick check: if you interview 10 people in the proposed segment, could the same initial product credibly serve at least several of them without extensive customization? If not, you may have selected a cluster of bespoke consulting opportunities rather than a market.

\subsection*{22. The Minimum Viable Market Is Not a Fixed Population Count}

There is no universally correct minimum of 10 customers, a thousand prospects, or ₹100 crore of TAM. A specialist enterprise firm might reach profitability with 4 annual clients. A low-priced consumer app may need tens of thousands of users. A high-touch school needs enough suitable learners to create pedagogical and economic viability within its cohort capacity.

The minimum useful segment is one large enough to sustain \textbf{the next stage of learning and operation} at an acceptable cost. It may be economically small but strategically excellent if it is reachable, serves an urgent job, and produces evidence transferable to a much larger adjacent market. It may be large but useless if winning the first transaction requires a distribution engine the founder does not possess.

\subsection*{23. Buyer Concentration and Customer Concentration Are Double-Edged}

If a beachhead contains 50 institutions that make large purchases, prospect identification is easy and contract values may be attractive. But sales may be slow, renewal may be lumpy, and losing 1 institution may erase a large portion of revenue.

If it contains 10,000 individuals making smaller purchases, diversification is better but acquisition may be expensive and transactions may be infrequent. A marketplace may have an even more complicated concentration problem if a handful of providers supply most of the transactions.

Identify the largest plausible revenue contribution of any 1 customer and what losing it would do to the business. A founding niche should not require 1 customer's goodwill to keep the organization solvent unless the arrangement is explicitly temporary and protected by reserves or other safeguards.

\subsection*{24. Test Whether the Segment Can Recognize Itself}

Good segmentation should make the intended customer say, ``That is my situation,'' without a lengthy explanation. It should also help the founder find people efficiently.

``For engineers who are assigned to evaluate an LLM application before release'' is identifiable through professional responsibilities. ``For visionary, high-agency lifelong learners who wish to unlock their potential'' is aspirational but operationally vague. The latter may serve as brand philosophy, not a practical acquisition definition.

The test is not whether the sentence is beautiful. It is whether the buyer can self-identify, outsiders can help locate eligible buyers, and the product's exclusions become intelligible.

\subsection*{25. Segment Economics Must Account for Expensive Outliers}

An apparently homogeneous group can contain customers whose service costs differ radically. One engineer may need 2 scheduled review sessions; another may require 20 hours of prerequisites. One college may have a cooperative department and timetable; another may require custom approvals, travel, and unpaid administrative work.

Measure not only the average cost but the \textbf{distribution of delivery load.} Set prerequisites, support limits, and acceptance conditions so that the same offer can be delivered with reasonably predictable economics. Otherwise the beachhead is a label on a set of unrelated jobs.

\subsection*{26. What Makes a Segment Strategically Reachable?}

Reachability has at least 4 components:

\begin{enumerate}
    \item The ability to identify individual prospective buyers;
    \item A lawful and ethical way to contact or be found by them;
    \item A credible route through their decision process; and
    \item The ability to produce enough transactions at sustainable acquisition cost.
\end{enumerate}

An elegant target such as ``all India-based researchers who might want advanced Bayesian modeling'' may be easy to define and hard to reach. Another group concentrated in a handful of well-run specialist professional communities may be easier to contact and qualify. A college market may be finite and identifiable, but purchase access can require institutional relationships and multistep approval.

The most useful early market is often one where \textbf{a few real interactions can be arranged immediately} rather than one that is impressive in a national industry report.

\subsection*{27. Segment Discovery as an Iterative Research Process}

Start with plausible groups from Chapter 5's opportunity map and Chapter 6's decision histories. Interview several members of each group, with deliberate inclusion of nonbuyers, existing users of alternatives, and buyers who recently switched. Extract common triggers, alternatives, budgets, and delivery constraints. Define candidate segments; challenge each with disconfirming cases; then test a narrowly priced offer.

Do not conduct segmentation as a single workshop after which the labels are frozen forever. Early paid customers may expose a much sharper pattern: perhaps engineers asked to evaluate production AI systems, not engineers who merely fear AI. The segment definition should become more accurate as real behavior appears.

The resulting segment is a theory. The first sale, successful delivery, repeat or referral, and subsequent independently acquired sales provide evidence that it may correspond to a sustainable market.
